\chapter{Conclusiones}
\label{cap:Conclusiones}

\section{Objetivos alcanzados}

Se ha alcanzado el objetivo principal del trabajo: el diseño e implementación de un prototipo funcional de codificación clínica automática en el estándar CIE-10-ES. El sistema resultante integra un clasificador multilabel basado en \emph{transformers} con una arquitectura completa de software que permite a un codificador clínico enviar informes de alta hospitalaria en español y recibir propuestas de codificación con su nivel de confianza asociado, sin necesidad de conocer los detalles del modelo subyacente.

A continuación se revisa el grado de cumplimiento de cada uno de los subobjetivos enumerados en el capítulo~\ref{cap:Objetivo}:

\begin{enumerate}
 \item \textbf{Adquisición y construcción de las fuentes de datos.} Completado. Se descargó el corpus CodiEsp en sus tres particiones (diagnósticos, procedimientos y variante extendida con anotaciones a nivel de \emph{span}), y se desarrollaron los scripts de \emph{scraping} sobre el portal eCIE-maps del Ministerio de Sanidad para obtener el catálogo completo de códigos CIE-10-ES, cuyo formato de descarga pública estructurada no existe. El resultado son los ficheros CSV que alimentan todas las fases posteriores del proyecto.

 \item \textbf{Módulo de preprocesamiento y gestión del corpus.} Completado. Se implementó el \emph{pipeline} de carga, limpieza y transformación de las anotaciones al formato multilabel requerido por el modelo, incluyendo la conversión de las anotaciones de nivel \emph{span} al nivel de documento empleado en el entrenamiento.

 \item \textbf{Análisis estadístico de las fuentes de datos.} Completado. El análisis de la distribución de longitudes, la frecuencia de aparición de los códigos y el desequilibrio de clases fundamentó las principales decisiones de diseño del modelo: el uso de ponderación de clases en la función de pérdida y la exploración de umbrales de decisión por código, finalmente descartados en favor de un umbral global (véase el subobjetivo siguiente).

 \item \textbf{Módulo clasificador multilabel.} Completado. Tras una evaluación comparativa de diez modelos candidatos (recogida en el Anexo~\ref{cap:AnexoA}), se seleccionó \texttt{IIC/RigoBERTa-Clinical} como modelo base y se ajustó mediante \emph{fine-tuning} sobre la partición de diagnósticos de CodiEsp, que es el alcance definido en el capítulo de objetivos; la codificación de procedimientos queda fuera del trabajo. El modelo final, entrenado sobre 500 documentos con 1.767 códigos objetivo, alcanza sobre el conjunto de prueba un \textbf{F1-micro de 0,489} y un \textbf{F1-macro de 0,109} (umbral global de 0,3, seleccionado en validación; los umbrales por clase se descartaron por sobreajustar el conjunto de validación). La calidad del \emph{ranking} de códigos, medida como \textbf{MAP por documento}, es de \textbf{0,437} sobre el gold completo, directamente comparable con el \emph{benchmark} CodiEsp-D 2020~\cite{miranda2020}, y de 0,537 al restringirla al vocabulario que el modelo puede predecir. El sistema se sitúa dentro del rango de los participantes del \emph{benchmark}, por detrás tanto de los \emph{ensembles} y diccionarios construidos manualmente que lo encabezan (IXA-AAA, MAP~0,593) como de los mejores sistemas basados exclusivamente en transformers (The Mental Strokers, MAP~0,517), lo que resulta razonable para un único clasificador plano entrenado con 500 documentos. La brecha entre F1-micro y F1-macro refleja el fuerte desequilibrio de clases del corpus, donde muchos códigos cuentan con muy pocos ejemplos de entrenamiento.

 \item \textbf{Backend de gestión y trazabilidad.} Completado. Se desarrolló un backend en Elixir con Phoenix que orquesta el flujo de análisis, gestiona la autenticación con confirmación por correo electrónico y aplica una arquitectura CQRS/\emph{Event Sourcing} sobre PostgreSQL, de modo que cada acción clínica queda registrada como un evento inmutable con su marca temporal y su autoría.

 \item \textbf{Servicio de inferencia (API REST).} Completado. El motor de IA expone los \emph{endpoints} \texttt{POST /predict} y \texttt{GET /} mediante FastAPI, con una imagen Docker optimizada para PyTorch y soporte para CPU y GPU. El servicio arranca incluso sin modelo entrenado, devolviendo un código 503 informativo, lo que facilita el desarrollo y las pruebas de integración.

 \item \textbf{Módulo de generación de resúmenes.} Completado. Se integró un modelo de lenguaje generativo ligero (cuantizado en formato GGUF, ejecutable en CPU) que produce un resumen o una paráfrasis del informe, con modelo, modo y \emph{prompts} configurables desde el \emph{backoffice} y emitido en tiempo real mediante \emph{streaming}. Por tratarse de una ayuda opcional y por su coste computacional, permanece desactivado por defecto en el despliegue de producción y se habilita bajo demanda. Su calidad se validó de forma cualitativa durante el desarrollo; una evaluación cuantitativa sistemática de los resúmenes generados (por ejemplo, con métricas de fidelidad factual y de cobertura de la información clínica) queda pendiente como trabajo futuro.

 \item \textbf{Interfaz web de demostración.} Completada. La interfaz web, desarrollada en Next.js 16 con React 19, permite enviar informes clínicos y visualizar en tiempo real las tarjetas de análisis generadas por el sistema. El codificador puede validar o rechazar individualmente cada código predicho y consultar el historial de análisis anteriores desde el panel lateral.
\end{enumerate}

Más allá de los subobjetivos enumerados, el desarrollo ha incluido componentes necesarios para alcanzar un prototipo funcional completo: el \emph{backoffice} de administración, la internacionalización en español e inglés, y una infraestructura de despliegue sobre VPS con Docker Compose y Traefik.

\section{Verificación de los requisitos}

La Tabla~\ref{tab:verificacion-req} resume el grado de cumplimiento de los requisitos definidos en el capítulo~\ref{cap:Objetivo}.

\begin{table}[h]
\centering
\small
\caption{Verificación de los requisitos funcionales (RF) y no funcionales (RNF).}
\label{tab:verificacion-req}
\begin{tabular}{llp{7.2cm}}
\hline
\textbf{Req.} & \textbf{Estado} & \textbf{Evidencia} \\
\hline
RF1 & Cumplido & El motor acepta texto clínico libre sin marcado previo. \\
RF2 & Cumplido & \texttt{POST /predict} devuelve los códigos de diagnóstico predichos. \\
RF3 & Cumplido & Cada código incluye su probabilidad y una confianza relativa normalizada. \\
RF4 & No cumplido & La inferencia por lotes no se implementó; el servicio procesa un informe por petición. \\
RF5 & Cumplido & API REST documentada con OpenAPI, con \emph{endpoint} de inferencia y de estado (Anexo~\ref{cap:AnexoM}). \\
\hline
RNF1 & Parcialmente & MAP 0,437 sobre prueba: dentro del rango del \emph{benchmark}, pero por debajo de los mejores sistemas (0,593 y 0,517). \\
RNF2 & Cumplido & Latencia inferior a 1\,s por informe en CPU de producción (Anexo~\ref{cap:AnexoI}). \\
RNF3 & Cumplido & Servicios desacoplados en contenedores con interfaces explícitas; el modelo base se sustituye sin tocar el resto. \\
RNF4 & Cumplido & \emph{Event Sourcing} registra cada acción con marca temporal, autoría y versión del modelo. \\
RNF5 & Cumplido & Exportación y borrado de cuenta, cifrado en tránsito y minimización de datos (sección~\ref{sec:rgpd}). \\
\hline
\end{tabular}
\end{table}

Los dos requisitos no satisfechos por completo se reconocen de forma explícita: la \textbf{inferencia por lotes} (RF4) quedó fuera del alcance final por no ser necesaria para el prototipo de demostración, y el \textbf{rendimiento} (RNF1) se cumple solo parcialmente, dado que el modelo se sitúa dentro del rango de los participantes del \emph{benchmark} pero no alcanza a los mejores sistemas. Ambas limitaciones se abordan en la sección de trabajos futuros.

\section{Justificación de competencias adquiridas}

En el TFG se han aplicado las competencias correspondientes a la Tecnología Específica de \emph{Computación}:

\begin{description}
 \item[CO4:] \emph{Capacidad para conocer y desarrollar técnicas de aprendizaje computacional y diseñar e implementar aplicaciones y sistemas que las utilicen, incluyendo las dedicadas a extracción automática de información y conocimiento a partir de grandes volúmenes de datos.} Esta competencia es el núcleo del trabajo: se ha diseñado, implementado y evaluado un clasificador multilabel de texto clínico basado en \emph{transformers}, realizando el ciclo completo de aprendizaje supervisado por transferencia (\emph{fine-tuning}) sobre el corpus CodiEsp. El proceso ha requerido la aplicación de técnicas de aprendizaje profundo: función de pérdida Binary Cross-Entropy con ponderación de clases para combatir el desequilibrio, planificación de la tasa de aprendizaje con \emph{cosine schedule} y \emph{plateau scheduler}, descongelación progresiva de capas del \emph{encoder}, y calibración de umbrales de decisión por clase sobre el conjunto de validación. El sistema resultante extrae automáticamente códigos CIE-10-ES a partir de informes clínicos en texto libre, y ha sido evaluado con métricas estándar de clasificación multilabel (F1-micro y F1-macro) sobre el \emph{benchmark} CodiEsp de referencia en español.
\end{description}

\section{Trabajos derivados y futuros}

Cinco líneas de trabajo concretas surgen del estado actual del prototipo.

\textbf{Ampliación del corpus de entrenamiento}. El clasificador se ha entrenado sobre 500 documentos de CodiEsp, lo que limita la cobertura de los códigos de enfermedades raras o procedimientos poco frecuentes, como refleja la diferencia entre el F1-micro de 0,489 y el F1-macro de 0,109 sobre el conjunto de prueba (umbral global de 0,3). A esta limitación se añade una restricción estructural del espacio de salida, inherente a los datos y no al modelo: el clasificador solo puede predecir los 1.767 códigos presentes en CodiEsp, en torno al 2,6\,\% de los 68.069 del catálogo CIE-10-ES completo. El resto de códigos no aparece en el corpus público: un clasificador supervisado plano no puede asignar un código ausente de su espacio de salida. Ampliar ese espacio exige ejemplos anotados o bien enfoques \emph{zero-shot} apoyados en las descripciones del catálogo, línea que este trabajo solo explora de forma parcial mediante el preentrenamiento débil. La vía para ampliar esa cobertura no es un esfuerzo puntual de anotación masiva, sino el ciclo de aprendizaje continuo descrito en la sección~\ref{sec:continuous-learning}: a medida que los codificadores validan o corrigen las predicciones en uso real, cada análisis aporta ejemplos etiquetados de nuevos códigos, de modo que el vocabulario que el modelo es capaz de reconocer crece de forma incremental con el propio uso del sistema. El F1-micro promedia ponderando por frecuencia de código, por lo que los códigos comunes dominan el resultado; el F1-macro da el mismo peso a cada código y penaliza directamente los raros con pocos ejemplos de entrenamiento. La incorporación de historias clínicas reales anonimizadas, en colaboración con centros hospitalarios, permitiría aumentar la representación de esos códigos y mejorar la utilidad clínica del sistema precisamente en los casos más atípicos. Esta colaboración representa el paso más directo para elevar el rendimiento del modelo.

\textbf{Migración a infraestructura cloud con soporte GPU}. El sistema se despliega actualmente sobre un VPS europeo con Docker Compose y Traefik, con un coste de $\sim$130\,€/año, ejecutando el motor de IA en modo CPU. Esta elección responde al coste: un nodo GPU en un proveedor cloud gestionado supera los 1.000\,€/mes, inasumible para un prototipo. Si el volumen de usuarios creciera y la latencia de inferencia pasara a ser un cuello de botella, la migración a un entorno con GPU (cloud gestionado, Kubernetes propio o instancia dedicada) no requeriría modificar ningún componente de la aplicación, al estar todo dockerizado. El Anexo~\ref{cap:AnexoI} recoge el detalle de la infraestructura actual y la comparativa de costes.

\textbf{Explorar la combinación de varias ejecuciones}. Al medir la variabilidad estocástica del entrenamiento se observó que promediar las salidas de tres ejecuciones disponibles reducía la varianza y mejoraba el MAP sobre el conjunto de prueba (sección~\ref{sec:anexoF-map}). Se trata de una observación incidental sobre ejecuciones que perseguían otros objetivos, no de un sistema diseñado ni evaluado como tal, por lo que no forma parte de los resultados de este trabajo. Diseñar y evaluar correctamente una combinación de modelos, midiendo su coste de inferencia frente a la ganancia, es una línea de trabajo con margen razonable de mejora. El obstáculo para llevarla a producción es precisamente ese coste: el motor de IA se ejecuta sobre un VPS sin GPU, donde multiplicar las pasadas del \emph{encoder} por informe no resulta viable. La vía que sortearía esa restricción es la \textbf{destilación de conocimiento}~\cite{hinton2015distilling}: entrenar un único clasificador contra las probabilidades producidas por la combinación, en lugar de contra las etiquetas binarias del corpus, conservando así el coste de inferencia de un solo modelo. Una prueba preliminar (sección~\ref{sec:anexoF-map}) muestra que esa vía no es inmediata en este régimen de datos: con 500 documentos de entrenamiento el profesor memoriza el corpus, sus probabilidades sobre él coinciden con las etiquetas binarias y el alumno no hereda nada del conjunto. Lo que queda abierto es destilar sobre un corpus clínico no etiquetado y ajeno al entrenamiento, donde el profesor sí produciría probabilidades informativas; requiere textos adicionales de los que este trabajo no dispone.

\textbf{Mejora de la explicabilidad}. El sistema atribuye la importancia de cada palabra por perturbación: sustituye la palabra por la máscara del \emph{tokenizador} y mide cuánto cae el \emph{logit} del código. Se adoptó tras descartar Gradient~$\times$~Input, que producía atribuciones ruidosas en una arquitectura de 24 capas por la difusión del gradiente (sección~\ref{sec:explicabilidad}). El método desplegado es fiel por construcción, porque mide el efecto real de quitar cada palabra, pero su coste crece con la longitud del informe y no cuantifica la incertidumbre de la atribución. Las dos extensiones naturales son \emph{SmoothGrad}~\cite{Smilkov17}, que promedia varias perturbaciones con ruido para estimar intervalos sobre la importancia, y \emph{SHAP KernelExplainer}~\cite{Lundberg17}, que reparte la importancia entre palabras con garantías de equidad. Ambas añaden coste por consulta, de modo que la pregunta que queda abierta es empírica: si la ganancia en calidad de la explicación justifica ese coste, medido sobre un conjunto de informes anotados con los términos clínicamente relevantes.

\textbf{Despliegue en entorno hospitalario real con captura de retroalimentación de los codificadores}. El prototipo incluye ya los mecanismos de validación y rechazo de códigos por parte del codificador, lo que sienta las bases para un ciclo de retroalimentación activa: las correcciones realizadas por los codificadores clínicos pueden usarse para reetiquetar documentos y realimentar el entrenamiento del modelo de forma iterativa. La puesta en marcha de este ciclo en un entorno real permitiría evaluar la utilidad del sistema en condiciones de uso auténticas y obtener datos de entrenamiento de mayor calidad que los disponibles en el corpus público.

\section{Viabilidad económica}

Más allá de su valor como prototipo académico, la arquitectura del sistema y su bajo coste operativo permiten esbozar un modelo de explotación económicamente viable, que se resume en los puntos siguientes.

\textbf{Modelo de suscripción por usuario}. La plataforma puede monetizarse mediante una suscripción mensual por usuario. Dado que el coste de alojamiento es actualmente de $\sim$130\,€/año, el umbral de rentabilidad es extremadamente bajo, lo que permite fijar una tarifa inicial asequible para facilitar la adopción. A medida que el número de usuarios crezca, la tarifa puede aumentarse de forma progresiva o diversificarse en planes con distintos niveles de uso.

\textbf{Planes de uso y límites de peticiones}. Una estructura de planes escalonados es la forma natural de monetizar el acceso a la API de predicción. El plan básico puede incluir un número limitado de análisis mensuales, mientras que los planes superiores ofrecen mayor volumen, acceso prioritario o integración directa con sistemas hospitalarios mediante API con límites de peticiones más amplios. Este modelo es habitual en servicios de IA como servicio (\emph{AI-as-a-Service}) y permite ajustar el precio al valor generado: un hospital con varios codificadores a tiempo completo genera un volumen de análisis muy grande, pudiendo incluso instalar instancias on-premise con soporte dedicado.

\textbf{Realimentación de datos: más usuarios, mejor modelo}. El ciclo de retroalimentación descrito en la sección~\ref{sec:continuous-learning} tiene además una lectura económica: cada usuario que valida o rechaza códigos aporta datos etiquetados por un profesional clínico, el tipo de dato más escaso y de mayor impacto sobre el modelo, de modo que un mayor volumen de usuarios alimenta ciclos de reentrenamiento que mejoran el sistema y refuerzan su utilidad.

\textbf{Escalado de infraestructura bajo demanda}. El coste de infraestructura crece solo cuando el volumen lo justifica: el VPS actual basta en la fase inicial y la migración a cloud o GPU descrita en los trabajos futuros no requiere reingeniería de la aplicación. La inversión en infraestructura es así proporcional al crecimiento, no un coste fijo de entrada.

\section{Valoración personal}

El proyecto ha supuesto un recorrido completo por el ciclo de vida de un sistema de \emph{software} con componente de investigación aplicada: desde la adquisición y el análisis de los datos hasta el despliegue en producción, pasando por la selección y el entrenamiento del modelo, el diseño de la arquitectura y el desarrollo de todos los componentes del sistema.

La parte más exigente y con mayor desviación respecto a la planificación original ha sido el ciclo de experimentación del modelo. Los primeros experimentos con los modelos candidatos no alcanzaban el nivel de rendimiento esperado, lo que obligó a repetir el ciclo de entrenamiento con distintas combinaciones de hiperparámetros, función de pérdida y estrategia de descongelación antes de dar por cerrada esa épica. Esta experiencia ha permitido desarrollar una intuición práctica sobre el impacto de cada decisión de diseño, tasa de aprendizaje, ponderación de clases, umbral de decisión, que difícilmente se adquiere de otro modo que con la experimentación directa y la observación de las curvas de aprendizaje.

Durante el desarrollo del código se empleó GitHub Copilot como asistente de autocompletado, agilizando la escritura de bloques repetitivos y la generación de esqueletos de código, y como herramienta de revisión de código mediante los agentes definidos en el repositorio del proyecto; en ambos casos, toda sugerencia fue revisada y adaptada al contexto del proyecto. Durante la redacción de la memoria se emplearon \texttt{aspell} para la corrección ortográfica y el complemento Grammarly para la corrección gramatical y de estilo.

El proyecto también ha puesto de manifiesto el valor de mantener los componentes del sistema con interfaces bien definidas. El uso del \emph{mock} del motor de IA para el desarrollo del backend y el frontend, con independencia del estado del modelo entrenado, ha evitado bloqueos entre las distintas líneas de trabajo y ha permitido avanzar en paralelo de forma efectiva, reduciendo el impacto de las iteraciones del ciclo de entrenamiento sobre el resto del sistema.

En términos cuantitativos, el proyecto se ha desarrollado a lo largo de aproximadamente 376 horas de trabajo efectivo, repartidas en 72 puntos de historia entre las seis épicas descritas en el capítulo~\ref{cap:Planificacion} y otros 22 dedicados a la redacción de la presente memoria, 94 puntos en total. El historial de desarrollo queda registrado en el repositorio público del proyecto~\cite{cie10repo2026}. Los primeros experimentos de entrenamiento se iniciaron en noviembre de 2024 en la rama \texttt{master}, acumulando conocimiento progresivo sobre el dominio, los datos y el comportamiento de los modelos candidatos. En enero de 2026, con todo ese aprendizaje consolidado y aprovechando los scripts, análisis y conclusiones ya obtenidos, se tomó la decisión de reescribir el historial Git y comenzar el desarrollo del sistema completo desde una base limpia, lo que permitió avanzar con mayor rapidez y solidez. La rama \texttt{main} resultante concentra más de 50 \emph{commits} del desarrollo definitivo del sistema.
